% !TEX program = pdflatex
% Example: Springer LNCS conference paper template
% Compile with: pdflatex main.tex && bibtex main && pdflatex main.tex && pdflatex main.tex
% Requires: llncs class (tlmgr install llncs)

\documentclass{llncs}

\usepackage{graphicx}
\usepackage{amsmath,amssymb}
\usepackage{booktabs}

% Self-contained bibliography
\usepackage{filecontents}
\begin{filecontents*}{references.bib}
@inproceedings{he2016,
  author    = {Kaiming He and Xiangyu Zhang and Shaoqing Ren and Jian Sun},
  title     = {Deep Residual Learning for Image Recognition},
  booktitle = {Proc. CVPR},
  year      = {2016},
  pages     = {770--778}
}
@article{hochreiter1997,
  author  = {Sepp Hochreiter and J\"urgen Schmidhuber},
  title   = {Long Short-Term Memory},
  journal = {Neural Computation},
  volume  = {9},
  number  = {8},
  pages   = {1735--1780},
  year    = {1997}
}
\end{filecontents*}

\begin{document}

\title{Efficient Neural Networks for Real-Time Signal Processing}
\author{Author Name}
\institute{Department of Electrical Engineering, University Name, \\
\email{author@university.edu}}

\maketitle

\begin{abstract}
This paper presents an efficient neural network architecture
for real-time signal processing applications. The proposed
model achieves competitive accuracy with significantly
reduced computational cost.

\keywords{Neural networks \and Signal processing \and Real-time systems}
\end{abstract}

\section{Introduction}
Real-time signal processing requires models that balance
accuracy with computational efficiency. Traditional deep
learning models~\cite{hochreiter1997} are often too heavy
for real-time deployment. Recent architectures~\cite{he2016}
have shown that residual connections can improve efficiency.

\section{Proposed Architecture}
We propose a lightweight architecture using depthwise
separable convolutions:
\begin{equation}
y = \text{DWConv}(x, k) \ast \text{PWConv}(W)
\end{equation}
where $k$ is the kernel size and $W$ are learned weights.

\section{Experiments}
Table~\ref{tab:results} compares our model with baselines.

\begin{table}[htbp]
\caption{Model comparison.}
\label{tab:results}
\centering
\begin{tabular}{lcc}
\toprule
Model & Accuracy & Parameters \\
\midrule
Full CNN & 94.2\% & 2.3M \\
MobileNet & 92.1\% & 0.5M \\
Proposed & 93.5\% & 0.3M \\
\bottomrule
\end{tabular}
\end{table}

\section{Conclusion}
Our lightweight architecture achieves 93.5\% accuracy with
only 0.3M parameters, making it suitable for real-time
deployment on edge devices.

\bibliographystyle{splncs04}
\bibliography{references}

\end{document}
